\chapter{Gestión del proyecto}

\section{Metodología}

El desarrollo se ha organizado mediante un enfoque iterativo e incremental. Las reuniones de seguimiento se celebraron de forma periódica, generalmente cada dos jueves, con participación de tutores, equipo técnico y personas vinculadas al proyecto original. En estas reuniones se revisaban avances, problemas, decisiones de arquitectura y siguientes objetivos. Esta dinámica se aproxima a una metodología ágil por sprints, aunque adaptada al contexto académico y a la disponibilidad del equipo.

La planificación inicial contemplaba una progresión amplia: semáforos y señales, IA vehicular, stops y cedas, perfiles de conductor, peatones, bicicletas, animales, diversidad de assets y optimización. El desarrollo real concentró el esfuerzo en construir una base sólida de IA vehicular, debido a la dificultad técnica de lograr navegación autónoma estable en la ciudad.

Durante el desarrollo se utilizaron asistentes de inteligencia artificial como apoyo de análisis y depuración. Su papel fue generar hipótesis, ordenar alternativas, proponer comprobaciones y ayudar a interpretar logs, Blueprints y CSV. Las respuestas no se incorporaron de forma automática: las decisiones se contrastaron con el comportamiento observado en Unreal Engine, las descripciones fechadas del programa, los resultados exportados y las pruebas de regresión. El anexo B documenta una selección depurada de estos prompts.

El seguimiento operativo se registró en GitLab mediante issues y milestones. El CSV depurado contiene 40 issues cerradas agrupadas en 15 sprints, desde el 7 de noviembre de 2025 hasta el 1 de julio de 2026. Esta fuente no sustituye a las actas de reunión ni a los resultados experimentales, pero sí aporta una trazabilidad intermedia: permite relacionar fechas, tareas cerradas y cambios técnicos concretos.

\Needspace{18\baselineskip}
\section{Seguimiento en GitLab}

La tabla~\ref{tab:milestones-gitlab} resume las milestones registradas en GitLab. Se han agrupado por sprint para evitar una lista excesiva de issues individuales y para conservar la lectura global de la evolución.

{\small
\setlength{\tabcolsep}{4pt}
\begin{longtable}{L{1.3cm}L{2.8cm}L{1.6cm}L{7.7cm}}
\caption{Milestones e issues cerradas registradas en GitLab.}\label{tab:milestones-gitlab}\\
\toprule
\textbf{Sprint} & \textbf{Periodo} & \textbf{Issues} & \textbf{Resultado documentado}\\
\midrule
\endfirsthead
\caption[]{Milestones e issues cerradas registradas en GitLab (continuación).}\\
\toprule
\textbf{Sprint} & \textbf{Periodo} & \textbf{Issues} & \textbf{Resultado documentado}\\
\midrule
\endhead
\bottomrule
\multicolumn{4}{>{\centering\arraybackslash}p{0.92\textwidth}}{\textit{\textbf{Fuente.} Elaboración propia a partir del CSV de issues y milestones de GitLab.}}\\
\endfoot
\bottomrule
\multicolumn{4}{>{\centering\arraybackslash}p{0.92\textwidth}}{\textit{\textbf{Fuente.} Elaboración propia a partir del CSV de issues y milestones de GitLab.}}\\
\endlastfoot
0 & 07/11 a 27/11/2025 & 2 & Reorganización inicial de \texttt{Content}, creación de \texttt{CarSimulator} y limpieza de assets heredados o importados.\\
1 & 27/11 a 11/12/2025 & 3 & Familiarización con la IA existente, semáforo dinámico y primeras señales verticales con zonas de detección y parada.\\
2 & 11/12 a 24/12/2025 & 3 & Primer vehículo autónomo con red inicial, mapeo de sensores y definición preliminar de perfiles de conducción con ASPAYM.\\
3 & 25/12/2025 a 08/01/2026 & 2 & Gestor evolutivo, persistencia de modelos y restricciones básicas para evitar estrategias de conducción no deseadas.\\
4 & 08/01 a 16/01/2026 & 2 & Primeros entrenamientos y tests centrados en corrección de desviación, junto con documentación de prácticas.\\
5 & 16/01 a 29/01/2026 & 2 & Ajuste de la red para mantener carril y entrenamiento inicial en curvas en ambos sentidos.\\
6 & 18/02 a 04/03/2026 & 4 & Bordes virtuales de intersección, mayor variedad de spawns, análisis de datos y multiraytracing.\\
7 & 05/03 a 19/03/2026 & 4 & Modo test, exportación ampliada de datos, normalización del fitness por spawn y mejora de puntos de aparición.\\
8 & 24/03 a 06/04/2026 & 4 & Integración de señales y semáforos en la red, zonas de espera, zonas de cruce y modo \textit{train/test}.\\
9 & 07/04 a 17/04/2026 & 3 & Telemetría ampliada, entradas \textit{one-hot} para contexto de tráfico y validación más segura de cruces.\\
10 & 20/04 a 30/04/2026 & 3 & Fitness avanzado, tiempo de gracia tras paradas legales, estadísticas de éxito y ampliación de la capa oculta.\\
11 & 01/05 a 15/05/2026 & 2 & Tasa de mutación dinámica, elitismo visual y refinamiento de la validación de paradas con tiempos aprendidos.\\
12 & 15/05 a 01/06/2026 & 2 & Penalización por colisión direccional, detección de \texttt{NAV\_VIOLATION} y ajuste de marcha atrás.\\
13 & 01/06 a 20/06/2026 & 2 & Lógica de rotondas, sensores de entrada y salida, conteo de giros y telemetría específica por salida.\\
14 & 22/06 a 01/07/2026 & 2 & Detección de solapes entre vehículos, distancia de seguridad y colas dinámicas de liberación en cedas.\\
\end{longtable}
}

\section{Cronología de decisiones}

La tabla~\ref{tab:cronologia} sintetiza los principales hitos recogidos en las actas de seguimiento.

\begin{longtable}{L{2.2cm}L{4.4cm}L{7.2cm}}
\caption{Principales hitos y decisiones documentadas en reuniones.}\label{tab:cronologia}\\
\toprule
\textbf{Fecha} & \textbf{Situación revisada} & \textbf{Decisión y consecuencia}\\
\midrule
\endfirsthead
\caption[]{Principales hitos y decisiones documentadas en reuniones (continuación).}\\
\toprule
\textbf{Fecha} & \textbf{Situación revisada} & \textbf{Decisión y consecuencia}\\
\midrule
\endhead
\bottomrule
\multicolumn{3}{c}{\textit{\textbf{Fuente.} Elaboración propia a partir de las actas de seguimiento.}}\\
\endfoot
\bottomrule
\multicolumn{3}{c}{\textit{\textbf{Fuente.} Elaboración propia a partir de las actas de seguimiento.}}\\
\endlastfoot
27/11/2025 & Integración de tráfico y elección de la primera IA. & Se prioriza la IA vehicular y se revisan semáforos dinámicos. Se descarta una malla dinámica para reconocer caminos por su coste de desarrollo.\\
11/12/2025 & Representación de las rutas de la ciudad. & Se descarta mapear todo el escenario con splines y se prueba el \textit{Path Finding} de Unreal. Se separan la decisión del modelo y la aplicación física del control.\\
17/12/2025 & Interfaz entre la red y el vehículo. & Se fijan dos salidas continuas, dirección y aceleración/frenado, ambas dentro de \texttt{[-1, 1]}.\\
08/01/2026 & Rigidez observada con \textit{Path Finding}. & Se abandona esta opción como solución principal y se adopta entrenamiento evolutivo por prueba y error para aprender las maniobras básicas.\\
05/02/2026 & Fallos en curvas e intersecciones. & Se introduce de forma provisional el etiquetado dinámico de bordes y se estudia cómo coordinar la elección de trayectoria con la percepción local.\\
19/02/2026 & Falta de una validación conjunta de los subsistemas. & Se programa una demostración con red neuronal, sensores y bordes activos en todo el mapa.\\
05/03/2026 & Demostración fallida y resultados difíciles de interpretar. & Se decide instrumentar la función de fitness y simplificar temporalmente los escenarios más complejos para localizar el origen de los fallos.\\
19/03/2026 & Segunda demostración con navegación ya estable en buena parte del mapa. & Se acepta la base técnica, se planifica la reorganización del repositorio y se preparan la compilación para Windows y las siguientes mejoras.\\
\end{longtable}

\Needspace{16\baselineskip}
\section{Línea evolutiva del desarrollo}

Además de los hitos de reunión, el proyecto ha seguido una evolución técnica marcada por pruebas que no siempre dieron un resultado utilizable a la primera. Esa evolución es relevante porque explica por qué la solución actual no se limita a una red neuronal aislada, sino que integra sensores, reglas de tráfico, persistencia de modelos, análisis de CSV y pruebas de regresión.

La tabla~\ref{tab:linea-evolutiva} resume las fases principales. No pretende listar cada cambio menor de Blueprint, sino recoger las decisiones que modificaron la dirección del trabajo.

{\small
\setlength{\tabcolsep}{4pt}
\begin{longtable}{L{2.2cm}L{4.2cm}L{7.5cm}}
\caption{Línea evolutiva del desarrollo técnico.}\label{tab:linea-evolutiva}\\
\toprule
\textbf{Periodo} & \textbf{Prueba o problema} & \textbf{Aprendizaje incorporado}\\
\midrule
\endfirsthead
\caption[]{Línea evolutiva del desarrollo técnico (continuación).}\\
\toprule
\textbf{Periodo} & \textbf{Prueba o problema} & \textbf{Aprendizaje incorporado}\\
\midrule
\endhead
\bottomrule
\multicolumn{3}{>{\centering\arraybackslash}p{0.9\textwidth}}{\textit{\textbf{Fuente.} Elaboración propia a partir de actas, issues de GitLab, descripciones de Blueprints y análisis de datos.}}\\
\endfoot
\bottomrule
\multicolumn{3}{>{\centering\arraybackslash}p{0.9\textwidth}}{\textit{\textbf{Fuente.} Elaboración propia a partir de actas, issues de GitLab, descripciones de Blueprints y análisis de datos.}}\\
\endlastfoot
Noviembre-diciembre & Tráfico autónomo inicial y conexión red-vehículo. & Se comprobó que el objetivo debía centrarse primero en IA vehicular. La red quedó acoplada a dos salidas continuas, dirección y aceleración/frenado, para respetar la física del vehículo.\\
Enero-febrero & Rigidez de NavMesh y \textit{Path Finding}. & Las rutas deterministas ayudaban a desplazar actores, pero no resolvían bien curvas, cruces ni control continuo. Se adoptó percepción local con sensores y aprendizaje evolutivo.\\
Marzo & Demostración fallida e interpretación insuficiente de los fallos. & La observación visual dejó de ser suficiente. Se instrumentó el fitness, se exportaron métricas por coche y se simplificaron temporalmente escenarios para localizar causas.\\
Marzo-abril & Precisión de parada, marcha atrás y coches inmóviles. & Se introdujeron penalizaciones específicas para parada imprecisa, retroceso y estancamiento injustificado. Las consultas de apoyo insistieron en evitar constantes arbitrarias y en derivar umbrales de variables ya presentes en el simulador.\\
Abril-mayo & Señalización, stops, cedas e intersecciones. & Los actores de tráfico pasaron a comunicar contexto al controlador. Los bordes de cruce se asociaron a dirección y trigger para que cada maniobra tuviera límites sensoriales coherentes.\\
Mayo-junio & Persistencia de modelos, duplicados y selección para reentrenar. & Se revisó la estrategia de \texttt{SuperCarModel}, \texttt{SessionCarModel} y copias fechadas. También se compararon respaldos para evitar repetir tests sobre modelos equivalentes y para elegir candidatos de reentrenamiento con datos, no solo por fitness máximo.\\
Junio-julio & Separación train/test, rotondas y análisis de regresión. & Se diferenciaron spawns y trayectorias por modo, se añadieron métricas de rotonda, solapes entre vehículos, bloqueo en cruce y colas dinámicas de liberación en cedas. La serie final muestra que una mejora local del fitness puede degradar otra parte del sistema, por lo que cada cambio debe validarse con bloques comparables antes de aceptarse.\\
\end{longtable}
}

\section{Alcance}

El alcance funcional completado se centra en conseguir que los vehículos de tráfico puedan aprender y circular dentro del escenario existente. El trabajo no modifica el vehículo conducido por el paciente ni rediseña la ciudad: actúa sobre los coches autónomos, su controlador y los elementos del mapa que les proporcionan información. El resultado es un controlador capaz de transformar las salidas de la red neuronal en dirección, aceleración y frenado sobre el sistema físico de Unreal Engine.

\Needspace{9\baselineskip}
Para alimentar ese controlador se ha incorporado percepción local mediante raycasting y se han normalizado las variables que entran en la red. El aprendizaje se realiza con poblaciones de vehículos, selección, mutación y elitismo. Los pesos obtenidos pueden guardarse y cargarse en sesiones posteriores, lo que evita reiniciar el proceso cada vez que se abre el simulador. La evaluación se apoya en una función de fitness que combina avance, centrado, velocidad, cumplimiento de señales y penalizaciones por comportamientos no deseados.

Incluye además semáforos, stops, cedas, líneas de parada y validación de maniobras. En cruces y rotondas se registran trayectorias, entrada, salida, giro, colisiones y recompensa, con puntos de aparición separados para entrenamiento y test. Los CSV resultantes permiten comparar ejecuciones y localizar fallos repetidos por zona.

No forman parte de esta versión la implementación completa de peatones, bicicletas y animales, ni los perfiles de conductor diferenciados previstos en la propuesta inicial. Tampoco se aborda todavía una validación con pacientes. Estas ampliaciones requieren primero consolidar la navegación vehicular y disponer de un protocolo de prueba estable, por lo que se mantienen como trabajo posterior.

\section{Gestión de riesgos}

La tabla~\ref{tab:riesgos} relaciona los riesgos más relevantes con su impacto, probabilidad y estrategia de mitigación.

{\small
\setlength{\tabcolsep}{4pt}
\begin{longtable}{L{3cm}L{3.6cm}L{1.8cm}L{5.1cm}}
\caption{Riesgos principales del proyecto.}\label{tab:riesgos}\\
\toprule
\textbf{Riesgo} & \textbf{Impacto} & \textbf{Prob.} & \textbf{Mitigación}\\
\midrule
\endfirsthead
\caption[]{Riesgos principales del proyecto (continuación).}\\
\toprule
\textbf{Riesgo} & \textbf{Impacto} & \textbf{Prob.} & \textbf{Mitigación}\\
\midrule
\endhead
\bottomrule
\multicolumn{4}{c}{\textit{\textbf{Fuente.} Elaboración propia a partir del seguimiento técnico del proyecto.}}\\
\endfoot
\bottomrule
\multicolumn{4}{c}{\textit{\textbf{Fuente.} Elaboración propia a partir del seguimiento técnico del proyecto.}}\\
\endlastfoot
Rigidez de navegación determinista & Comportamiento poco realista y baja utilidad clínica & Alta & Pivotar a percepción local con sensores y red neuronal.\\
Fitness mal calibrado & Aprendizaje de estrategias no deseadas & Alta & Instrumentar componentes, exportar CSV y analizar penalizaciones por familia de muerte.\\
Coste computacional & Entrenamientos lentos y dependencia de GPU/CPU & Media & Poblaciones moderadas, análisis offline y propuesta futura de modo sin renderizado.\\
Complejidad de intersecciones & Colisiones y violaciones normativas en cruces & Alta & Bordes dinámicos, zonas de cruce, validación de stops/cedas y separación train/test.\\
Sobreajuste a spawns concretos & Buen rendimiento local pero mala generalización & Media & Normalización por punto de aparición y test con spawns diferenciados.\\
Pérdida de conocimiento entrenado & Reinicio de aprendizaje entre sesiones & Media & Persistencia de modelos \texttt{SuperCarModel} y \texttt{SessionCarModel}.\\
\end{longtable}
}

\Needspace{16\baselineskip}
\section{Estimación de recursos y presupuesto}

El proyecto se ha desarrollado durante las prácticas y el TFG, sin contratación externa ni compra específica de material. El recurso principal ha sido el tiempo dedicado a implementar Blueprints, ejecutar entrenamientos, revisar resultados y documentar cambios. Parte de ese tiempo no es programación directa: los entrenamientos ocupan el equipo durante periodos prolongados y obligan a repetir pruebas cuando cambia el fitness o la geometría de una intersección.

El trabajo técnico se ha realizado en un equipo capaz de ejecutar Unreal Engine 5 y las físicas de varios vehículos de forma simultánea. El control de versiones se ha gestionado con GitLab y el análisis con Python y bibliotecas libres. Unreal Engine no ha generado coste de licencia en este contexto académico. Por tanto, no existen costes directos de software ni de infraestructura; el hardware ya estaba disponible y solo correspondería imputar amortización y consumo eléctrico si se elaborase un presupuesto comercial.

En la memoria final, la valoración económica deberá convertir las horas efectivas de desarrollo y experimentación a un coste de ingeniería acordado con el tutor, y añadir la parte proporcional del equipo utilizado. Mantener separadas ambas partidas evita presentar como gasto real una licencia no pagada y permite actualizar el presupuesto cuando se cierre el registro definitivo de dedicación.
